\section{Introduction}

The recorded lecture is provided \myhref{https://www.youtube.com/watch?v=LvucqRh7JKI}{here}.

\begin{ejercicio}
    What are the two views on IoT discussed in the lecture? How do they differ?\\

    The two views on IoT are:
\begin{itemize}
    \item IoT as a network of connected \emph{things}, aka \emph{M2M (Machine to Machine)}.
    \item IoT as a network of connected \emph{data}, aka \emph{Big Data}.
\end{itemize}

    They differ in plenty of different aspects:
\begin{itemize}
    \item While the first one focuses on the devices and the communication between them, the second one focuses on the data that is generated by these things and how it can be used to create value.
    \item While the first one sees Internet as just a communication medium, the second one sees Internet as a platform for data collection and analysis.
    \item While the main goal of the first one is to connect as many things as possible, the main goal of the second one is to extract insights from the data.
    \item While the first one is usually taken by engineers and computer scientists, the second one is usually taken by business people and data scientists.
\end{itemize}
\end{ejercicio}

\begin{ejercicio}
    Are IoT gateways always necessary? What are they used for?\\

    IoT gateways are used to connect IoT things to the internet. IoT things are usually resource-constrained devices that cannot connect directly to the internet, so they need a gateway to act as a bridge between them and the internet. However, IoT gateways are not always necessary, as some IoT things can connect directly to the internet using low-power communication technologies such as Wi-Fi or cellular. In this case, the IoT thing itself can act as a gateway, and there is no need for an additional device.
\end{ejercicio}